%%%%%%%%%%%%%%%%%%%%%%%%%%%%%%%%%%%%%%%%%%%%%%%%%%%%%%%%%%%%%%%%%%%%%%%%%%%%
%  OPT 2026 BODY -- a fork of aaai_article/signmuon_body.tex, cut to the
%  workshop's five pages.  This file is edited by hand; make_opt.py reads it
%  instead of the AAAI body and reports when the AAAI body has moved on (it
%  records its hash in opt_body.origin).
%
%  What was cut, and where it went.  Most of it already had a home in the
%  appendix, so most of the cut is a deletion and a pointer, not a move:
%
%    Related Work (3 paragraphs)        -> opt_appendix_extra.tex
%    the Sign A family and its display  -> opt_appendix_extra.tex
%    the federated formulation (2)      -> opt_appendix_extra.tex
%    Sec. 4.1, the centralized setting  -> App. \ref{app:lrscale} (already there)
%    the unit-gain justification        -> App. \ref{app:lrscale} (already there)
%    Sec. 4.5, federated mechanics      -> App. \ref{app:alg} (already there)
%    the EF21 estimator display (10)    -> App. \ref{app:hyp} (already there)
%    the two-model display (11)         -> opt_appendix_extra.tex
%    memory cost, the sqrt(r) penalty   -> App. \ref{app:hyp} (already there)
%    Sec. 5.2, federated experiments    -> opt_appendix_extra.tex, with Table 1
%    the CIFAR and nanoGPT tables       -> opt_appendix_extra.tex
%
%  The two experiment figures came the other way, out of the appendix and into
%  Secs. 4.1 and 4.2: at five pages a curve earns its space and the digits
%  behind it do not.  What was left was then cut by about a third again,
%  sentence by sentence.
%
%  Labels the appendix points at must survive somewhere, or its references go
%  undefined.  These are kept here deliberately, however tempting the cut:
%  as:1, as:2 (19 references), eq:three_placements, eq:lin_obj,
%  eq:lin_increment, eq:efsm_est, eq:unit_gain_main, th:ef_div,
%  fig:divergence_plot, fig:cifar_results, fig:nanogpt, sec:proState,
%  sec:theory, sec:exp_nanogpt.  eq:signa_general, eq:two_models,
%  tab:cifar_main and tab:nanogpt are kept in opt_appendix_extra.tex.
%  make_opt.py checks all of this before compiling.
%%%%%%%%%%%%%%%%%%%%%%%%%%%%%%%%%%%%%%%%%%%%%%%%%%%%%%%%%%%%%%%%%%%%%%%%%%%%

\maketitle

\begin{abstract}\label{sec:abs}
  \begin{claude}
  SignMuon compresses the Muon update to one bit per parameter by taking the elementwise sign of the Muon Linear Minimization Oracle (LMO) direction: the most direct way to give a matrix-aware optimizer a one-bit budget. It outperforms SignSGD empirically, yet it can ascend even on a linear function. Signing the gradient \emph{before} the oracle rather than after does not repair this: we exhibit a small explicit instance on which sign-before (MuonUSign) and sign-on-both-sides (MuonSign) ascend as well, so no placement of the sign around the oracle descends in general. Error feedback does not repair sign-after either: applied to the \emph{output} of the oracle it fails for every smoothness constant, step size and momentum. Applied to the \emph{gradient} it does work, and EF21-MuonUSign and EF21-MuonSign attain the standard $\mathcal{O}(T^{-1/2})$ rate for the squared gradient norm on smooth nonconvex problems, the latter at one bit in each direction. Our experiments then reverse this order: across CIFAR-10, its federated split and the nanoGPT speedrun the strongest compressed method is sign-\emph{after}, precisely the placement we prove divergent, with the provably convergent variants behind it.
  \end{claude}
\end{abstract}

\section{Introduction}\label{sec:intro}
In distributed training over a limited link, communication rather than computation is the binding cost. Compression is the standard remedy \cite{bernstein2018signsgd,bernstein2019signsgdmajority,beznosikov2020biased}. SignSGD is its extreme point: one bit per coordinate, at a small cost in accuracy, but it discards matrix structure. Muon exploits that structure, orthonormalizing the momentum matrix before stepping \cite{jordan2024muon,bernstein2024old,shah2025practical}\cl{: its step is the Linear Minimization Oracle (LMO) of the spectral norm, and it is transmitted as a \emph{dense} matrix at full precision. We compress that step with a sign, placed \emph{after} the oracle, \emph{before} it, or on both sides:} \textbf{SignMuon}, $\operatorname{sign}(\operatorname{LMO}(\cdot))$; \textbf{MuonUSign}, $\operatorname{LMO}(\operatorname{sign}(\cdot))$; and \textbf{MuonSign}, $\operatorname{sign}(\operatorname{LMO}(\operatorname{sign}(\cdot)))$. SignMuon and MuonSign emit a $\pm1$-valued step, so uplink and downlink alike cost one bit.

\cl{SignMuon, MuonUSign and MuonSign do not descend in general: each} can \emph{ascend} on a linear objective at every step size and every momentum, SignMuon on a $4\times4$ gradient (Theorem~\ref{th:1}), MuonUSign and MuonSign on \cl{a $5\times5$ gradient they share} (Theorems~\ref{th:2}--\ref{th:3}). Error feedback is the standard repair for a biased compressor. \cl{For sign-after it does not help}: applied to the \emph{output} of the oracle it fails for every triple $(L,\eta,\mu)$ (Theorem~\ref{th:ef_div}). No step-size rule built from the smoothness and momentum constants can make it convergent. Applied to the \emph{gradient}, however, it restores convergence: adapting EF21-Muon~\citep{gruntkowska2025error} gives \textbf{EF21-MuonUSign} and \textbf{EF21-MuonSign}, which \cl{attain $\min_{t\le T}\mathbb{E}\|\nabla f(\mathbf{X}_t)\|_*^2=\mathcal{O}(T^{-1/2})$, the standard rate for smooth nonconvex problems}, at a one- and a two-sided one-bit budget.

\cl{The guarantees do not predict which method performs best in practice.} On centralized CIFAR-10 with a ResNet-18, on federated CIFAR-10 (Appendix~\ref{app:opt:fed}), and on the nanoGPT speedrun, the \cl{strongest} compressed method \cl{signs \emph{after} the oracle}, the placement error feedback cannot repair\cl{, and} EF21-MuonUSign and EF21-MuonSign consistently trail it. \cl{The counterexamples are worst-case constructions; random instances do not reproduce them (Section~\ref{sec:theory}).}

\paragraph{Related work.} We build on the interpretation of Muon as an LMO method \citep{bernstein2024old,pethick2025training,riabinin2025gluon} and on the error-feedback framework of \citet{gruntkowska2025error} for bidirectionally compressed Muon. The sign-before and both-sides placements are new, as are the two EF21 methods; SignMuon was proposed concurrently by \citet{mishra2026signmuon}, whose guarantee, we show in Appendix~\ref{app:mishra}, applies to a different method. Appendix~\ref{app:opt:related} gives the full account.

\section{Problem Statement}\label{sec:proState}

We consider $\min_{\mathbf{X} \in \mathcal{X}} \{f(\mathbf{X}) := \mathbb{E}_{\xi \sim \mathcal{D}} [f(\mathbf{X};\xi)]\}$ with $f(\cdot\,;\xi)$ continuously differentiable and possibly non-convex; the federated form, over $N$ clients, is \eqref{eq:opt:fed} in Appendix~\ref{app:opt:fed}. The model is a tuple of layers $\mathbf{X} = [\mathbf{X}_1,\dots,\mathbf{X}_p]$, $\mathbf{X}_i \in \mathbb{R}^{m_i\times n_i}$, with $\nabla_i$ the gradient block of layer $i$. This is the Gluon parameter space \citep{riabinin2025gluon} under a unified geometry: each layer is equipped with the spectral norm $\|\cdot\|_{2\to2}$ with nuclear dual $\|\cdot\|_{*}$, and $\|\nabla f(\mathbf{X})\|_*^2 := \sum_i\|\nabla_i f(\mathbf{X})\|_*^2$. On a block of width one the oracle only normalizes, so all three methods coincide with SignSGD; the results below need $\min(m_i,n_i)\ge2$ (Appendix~\ref{app:vector}).

\begin{assumption}[\textbf{Lower boundedness}]\label{as:1}
$f(\mathbf{X}) \ge f^*$ for all $\mathbf{X}$; where explicitly invoked, each local $f_j \ge f_j^*$ as well.
\end{assumption}

\begin{assumption}[\textbf{Layer-wise smoothness}]\label{as:2}
For every layer $i$ and all $\mathbf{X},\mathbf{Y}$, $\|\nabla_i g(\mathbf{X}) - \nabla_i g(\mathbf{Y})\|_{*} \le L_i^{g}\,\|\mathbf{X}_i - \mathbf{Y}_i\|_{2\to 2}$ for $g = f$ (constant $L_i$) and $g = f_j$ (constant $L_{i,j}$). Appendix~\ref{app:prelim} gives a weaker $(L^0,L^1)$ form, used only in Corollary~\ref{cor:l0l1}, and records the strength of the quantifier over arbitrary $\mathbf{X},\mathbf{Y}$.
\end{assumption}

An LMO minimizes the first-order model over the unit ball of a norm, returning the steepest-descent direction in it. For $\mathbf{M}\in\mathbb{R}^{m\times n}$ with rank-$r$ SVD $\mathbf{U}\boldsymbol{\Sigma}\mathbf{V}^\top$ the LMO direction is $A(\mathbf{M})=-\mathbf{U}\mathbf{V}^\top$; we write $\operatorname{polar}(\mathbf{M}):=-A(\mathbf{M})$, an orthogonalization, unscaled. The reference implementation and the RMS$\to$RMS interpretation rescale it by positive per-layer constants, invisible to $\operatorname{sign}(\cdot)$\cl{; we absorb them into the step size (Appendices~\ref{app:prelim} and~\ref{app:lrscale}).}

\section{Theory}\label{sec:theory}

\paragraph{Three placements of the sign.} All three build the same momentum estimate $\tilde{\mathbf{M}}_t$ from the stochastic gradients: an exponential moving average at coefficient $\mu\in[0,1)$, or its Nesterov variant (Appendix~\ref{app:opt:signa}). All three then step $\mathbf{X}_t=\mathbf{X}_{t-1}-\eta_t\,\mathbf{s}_t$ at learning rate $\eta_t$. They differ only in where the sign acts on $\tilde{\mathbf{M}}_t$:
\begin{equation}
\label{eq:three_placements}
\mathbf{s}_t=
\begin{cases}
\operatorname{sign}\!\bigl(\operatorname{polar}(\tilde{\mathbf{M}}_t)\bigr) & \text{(SignMuon: after)},\\[2pt]
\operatorname{polar}\!\bigl(\operatorname{sign}(\tilde{\mathbf{M}}_t)\bigr) & \text{(MuonUSign: before)},\\[2pt]
\operatorname{sign}\!\bigl(\operatorname{polar}(\operatorname{sign}(\tilde{\mathbf{M}}_t))\bigr) & \text{(MuonSign: both)}.
\end{cases}
\end{equation}
Throughout, $\operatorname{sign}(0)$ is resolved to an independent random $\pm1$, keeping the transmitted alphabet binary. \cl{Every method transmits one bit per matrix parameter on the uplink. On the downlink SignMuon and MuonSign distribute a $\pm1$-valued matrix and cost one bit again, whereas MuonUSign applies $\operatorname{polar}(\cdot)$ on the server, which is dense and goes at full precision.} Appendix~\ref{app:alg} gives the algorithms and the federated protocol.

% Declared ahead of its first reference: a float can only move forward on the
% page, and from the next subsection it would land after the text discussing it.
\begin{figure}[!t]
    \centering
    \includegraphics[width=0.86\textwidth]{images/counterexamples/counterexamples_main.pdf}
    \caption{All eight methods on the three counterexample instances; the ascending method is drawn heavy and named in each panel. \emph{Left:} SignMuon, $4\times4$ (Theorem~\ref{th:1}). \emph{Centre:} MuonUSign and MuonSign on their shared $5\times5$ instance (Theorems~\ref{th:2}--\ref{th:3}); both panels plot~\eqref{eq:lin_obj}. \emph{Right:} EF21-SignMuon, $2\times2$ (Theorem~\ref{th:ef_div}), normalized units at $\eta=1$. Trajectories are momentum-free without loss (Proposition~\ref{prop:reduction}; Lemma~\ref{lem:efsm_realize} for EF21-SignMuon).}
    \label{fig:divergence_plot}
\end{figure}

\subsection{Divergence of the Three Sign Placements}
The sign step is itself an LMO, steepest descent in $\ell_\infty$ \citep{bernstein2024old}. Signing a Muon direction therefore composes two oracles, each sound alone. The composition is an LMO for \emph{no} norm. An oracle for a norm returns $\mathbf{D}=-A(\mathbf{G})$ with $\langle\mathbf{G},\mathbf{D}\rangle=\|\mathbf{G}\|_{\mathrm{dual}}\ge0$; Theorem~\ref{th:1} provides a gradient on which the composition drives that inner product strictly negative.

We refute the descent property on \emph{linear} objectives:
\begin{equation}
\label{eq:lin_obj}
f(\mathbf{X})=\langle \mathbf{G},\mathbf{X}\rangle,\qquad \nabla f(\mathbf{X})\equiv \mathbf{G}.
\end{equation}
Gradient descent and full-precision Muon drive $f\to-\infty$ here, so a rule that drives $f\to+\infty$ is unambiguously ascending. (Remark~\ref{rem:efsm_bdd} restores Assumption~\ref{as:1} without moving any trajectory.) Whatever the momentum, the step is the constant matrix $\mathbf{s}(\mathbf{G})$, so momentum affects neither convergence nor divergence (Proposition~\ref{prop:reduction}). The increment is
\begin{equation}
\label{eq:lin_increment}
f(\mathbf{X}_{t})-f(\mathbf{X}_{t-1})=-\eta_t\,\bigl\langle \mathbf{G},\ \mathbf{s}(\mathbf{G})\bigr\rangle .
\end{equation}
A single $\mathbf{G}$ with $\langle\mathbf{G},\mathbf{s}(\mathbf{G})\rangle<0$ therefore makes $f$ increase at every step. Such a $\mathbf{G}$ exists for each placement at small size.

\begin{summarythm}{Theorems~\ref{th:1}--\ref{th:3} (summary)}
There exists $\mathbf{G}\in\mathbb{R}^{4\times4}$ with $\langle\mathbf{G},\mathbf{s}(\mathbf{G})\rangle<0$ for SignMuon, and $\mathbf{G}\in\mathbb{R}^{5\times5}$ with $\langle\mathbf{G},\mathbf{s}(\mathbf{G})\rangle<0$ for MuonUSign and for MuonSign simultaneously. On the corresponding objective~\eqref{eq:lin_obj} each method ascends at every iteration, for every $\eta_t>0$, every $\mu\in[0,1)$ and either momentum rule, and $f(\mathbf{X}_t)\to+\infty$ whenever $\sum_t\eta_t=\infty$.
\end{summarythm}

Both instances are explicit, both inner products exact rationals. Appendices~\ref{app:proof_th1}--\ref{app:proof_th3} state the theorems and prove them; Figure~\ref{fig:divergence_plot} (left, centre) shows the ascents. Negative results of this kind exist for uncompressed Muon as well \citep{parshakova2026muondoesconvergeconvex}. Ours are due to the compressor rather than the geometry.

%The counterexamples bound the theory from below in another sense as well. 
On a deterministic convex quadratic the alignment $\rho_t=\langle\nabla f,\mathbf{D}_t\rangle/(\|\nabla f\|_F\|\mathbf{D}_t\|_F)$ between the gradient and the update direction actually taken can be measured directly. Theorems~\ref{th:1}--\ref{th:3} give instances with $\rho_t<0$; on random instances all three placements keep $\rho_t$ away from zero (Appendix~\ref{app:images_task}). The constructions therefore describe a worst case rather than a typical one. The experiments below run contrary to them.

\subsection{Error Feedback on the Polar Factor Fails}

The classical remedy for a biased compressor is \emph{error feedback}: keep the discarded residual and fold it into the next message \citep{seide20141bit,karimireddy2019efsign}. We work in its EF21 form \citep{richtarik2021ef21}, which stores an estimator and compresses the \emph{difference} to it. No bounded-gradient assumption is then needed. The most direct use keeps the geometry, signing \emph{after} the LMO, and applies error feedback to the polar factor. \textbf{EF21-SignMuon} (Algorithm~\ref{ef21_signmuon}) tracks an estimate $\mathbf{d}_t^{\mathrm{est}}$ of the polar factor $\mathbf{D}_t=\operatorname{polar}(\tilde{\mathbf{M}}_t)$ by a scaled sign of the residual,
\begin{equation}\label{eq:efsm_est}
\mathbf{d}_t^{\mathrm{est}}=\mathbf{d}_{t-1}^{\mathrm{est}}+\alpha_t\operatorname{sign}\!\bigl(\mathbf{D}_t-\mathbf{d}_{t-1}^{\mathrm{est}}\bigr),
\qquad
\alpha_t=\operatorname{mean}\bigl|\mathbf{D}_t-\mathbf{d}_{t-1}^{\mathrm{est}}\bigr|,
\end{equation}
then steps $\mathbf{X}_t=\mathbf{X}_{t-1}-\eta_t\,\mathbf{d}_t^{\mathrm{est}}$. This does not repair the placement.
% The three divergence theorems are stated in the appendix, next to their
% proofs, but the paper's numbering is fixed by everything that cites it (the
% code, its READMEs, the reproduction guide): Theorems 1-3 are the three sign
% placements, 4 is this one, 5 is the convergence result. The counter is set by
% hand here and reset in the appendix to keep those numbers.
\setcounter{theorem}{3}
\begin{theorem}[Divergence of EF21-SignMuon]\label{th:ef_div}
For every $L>0$, step size $\eta>0$, momentum coefficient $\mu\in[0,1)$, and either momentum variant, there is an $L$-smooth (Assumption~\ref{as:2}), bounded-below (Assumption~\ref{as:1}) $f:\mathbb{R}^{2\times2}\!\to\mathbb{R}$ on which EF21-SignMuon started at $\mathbf{X}_0=\mathbf{0}$ diverges: for an explicit $c=c(f)>0$, $f(\mathbf{X}_{t+2})-f(\mathbf{X}_t)=c\,L\eta^2>0$ for $t\ge3$, so $f(\mathbf{X}_t)\to+\infty$. No step-size rule $\eta=\eta(L,\mu)$ using only the smoothness and momentum constants can make the method convergent.
\end{theorem}
The construction (Appendix~\ref{app:proof_th4}) turns the shared magnitude against the method. A $2\times2$ target whose large off-diagonal reverses sign every step holds $\alpha_t$ at $\Theta(1)$; the small diagonal constant is then overshot at every step, its estimate settling into a period-two cycle whose average has the \emph{wrong} sign. Neither dimension nor momentum is implicated (Appendix~\ref{app:proof_th4}): $2\times2$ is where the worst-case contraction $1/d$ of the scaled sign is most favourable.

\subsection{Error Feedback on the Gradient Works}
The defect is the target, not the shared magnitude. Error feedback needs a target that varies with the step size; $\operatorname{polar}(\tilde{\mathbf{M}}_t)$ is not Lipschitz in its argument, so it can move by $\Theta(1)$ between consecutive rounds however small $\eta$ is. Compressing the \emph{gradient} in its place restores that dependence. \cl{The resulting method is EF21-Muon~\citep{gruntkowska2025error} with the scaled sign $\mathcal{C}(\mathbf{Y}) = \operatorname{mean}|\mathbf{Y}|\,\operatorname{sign}(\mathbf{Y})$ on every compressed link: one bit per parameter, plus one scalar per layer.} \textbf{EF21-MuonUSign} takes the pair $(\mathcal{C},I)$: a compressed uplink, a full-precision model back. It maintains a gradient estimator $\mathbf{g}_t^{\mathrm{est}}$ from the residual $\tilde{\mathbf{M}}_t-\mathbf{g}_{t-1}^{\mathrm{est}}$ exactly as in \eqref{eq:efsm_est}, then steps $\mathbf{X}_t=\mathbf{X}_{t-1}-\eta_t\mathbf{D}_t$ with $\mathbf{D}_t=-A(\mathbf{g}_{t}^{\mathrm{est}})$, the polar factor of the estimator rather than of $\nabla f$. \textbf{EF21-MuonSign} takes $(\mathcal{C},\mathcal{C})$, one bit each way. Its downlink is a second error-feedback loop, on the \emph{model} rather than the gradient\cl{, splitting the method into a server model $\mathbf{X}_t$, which the guarantee bounds, and a broadcast model $\mathbf{W}_t$, where every gradient is evaluated (Appendix~\ref{app:opt:fed}).}

Both descend where the placements they repair ascend: Figure~\ref{fig:divergence_plot} (centre) has them on the $5\times5$ instance. Appendix~\ref{app:hyp} shows them to be \emph{exact instances} of the EF21-Muon framework; the one substantive check is that the scaled sign is contractive (Lemma~\ref{lem:signcontr})\cl{, so they inherit its guarantees}. Under Assumptions~\ref{as:1}--\ref{as:2} both reach $\min_{t\le T}\mathbb{E}\|\nabla f(\mathbf{X}_t)\|_*^2=\mathcal{O}(T^{-1/2})$, the rate uncompressed Muon attains: one bit changes the constant, not the rate. The second loop costs a step size smaller by \cl{a factor $\sqrt{r}$ in the layer rank. The factor arises because the scaled sign is contractive in the Euclidean norm and not in the layer norm the spectral geometry requires (Remark~\ref{rem:normmismatch}).}

\section{Experiments}\label{sec:exp}

We compare all six sign-based methods against Muon, SignSGD, SGD and Adam. Only the base rate $\eta_0$ is tuned: layer $\ell$, of shape $m\times n$, gets $\eta_\ell=\eta_0\lambda_\ell$ with
\begin{equation}
\label{eq:unit_gain_main}
\lambda_\ell=\sqrt{m}\,/\,\|\mathbf{P}_\ell\|_F
\quad\Longrightarrow\quad
\lambda_\ell^{\textsc{lmo}}=\sqrt{\max(1,m/n)},
\qquad
\lambda_\ell^{\textsc{sign}}=1/\sqrt{n},
\end{equation}
where $\mathbf{P}_\ell$ is the matrix the method applies, whose Frobenius norm is known exactly in either family. The criterion is that every layer receive the same root-mean-square input--output gain \citep{yang2023spectral,large2024modular}. Its \textsc{lmo} branch reproduces the aspect-ratio factor the reference Muon implementation applies \citep{jordan2024muon}\cl{; we apply} the same criterion \cl{to} the \textsc{sign} family, for which no convention exists (Appendix~\ref{app:lrscale}). 

\subsection{Image Classification}
We train a ResNet-18 adapted to low-resolution images on CIFAR-10 \cite{krizhevsky2009cifar} for $75$ epochs, $\eta_0$ cosine-annealed to zero and selected per method on a held-out split. Numbers average three seeds; differences below the seed spread are not claimed (Appendix~\ref{app:repro}).

\begin{figure}[!tb]
    \centering
    \includegraphics[width=0.91\textwidth]{images/centralized_images/cifar_main.pdf}
    \caption{Centralized ResNet-18 on CIFAR-10: test accuracy at the $\eta_0$ selected for each method, from epoch $25$ (Table~\ref{tab:cifar_main}). Panels group the sign placements, the EF21 variants and the baselines; Muon is the gray dashed reference, bands $\pm1$ s.d.\ over three seeds. The numbers are Table~\ref{tab:cifar_main}, Appendix~\ref{app:opt:cen}.}
    \label{fig:cifar_results}
\end{figure}


SignMuon ranks first, by less than one seed spread over Muon (Figure~\ref{fig:cifar_results}). The separation that resolves is the $1.2$ points over SignSGD. SignSGD spends the same budget without the \textsc{lmo}, so the geometry accounts for the gap, not the compressor. Error feedback is not free here: EF21-SignMuon lies $0.29$ points below SignMuon, against standard deviations of $0.11$ and $0.15$. Both families order the placements after, before, both sides. 
%Federation sharpens the ordering (Appendix~\ref{app:opt:fed}): SignMuon exceeds SignSGD by $4.3$ points there, at one bit in both directions. 

\subsection{Language Modelling}\label{sec:exp_nanogpt}

\begin{figure}[!tb]
\centering
\includegraphics[width=0.43\textwidth]{images/nanogpt/fig_nanogpt_main.pdf}
\caption{NanoGPT speedrun, $8\times$H100: validation loss against optimizer step, one run per method; dotted is the target $3.28$, and EF21-MuonSign is drawn at its broadcast model $\mathbf{W}$. The inset magnifies the boxed tail, where SignSGD lies above the range shown.}
\label{fig:nanogpt}
\end{figure}


We test where the matrices are large enough for the layer-rank dependence of Corollary~\ref{cor:smooth} to take effect: \cl{the modded-nanoGPT speedrun (record \#40) on $8\times$H100 \citep{jordan2024moddednanogpt}, whose hidden matrices are $768\times3072$. Each method replaces Muon on the hidden matrices and gates, at rates fixed a priori by \eqref{eq:unit_gain_main}, one $\eta_0$ per \emph{family}, so every contrast is matched-hyperparameter (Appendix~\ref{app:nanogpt}).}

The composition transfers to language modelling. All six such methods improve on SignSGD by at least $0.08$ in validation loss; the best two, EF21-SignMuon and SignMuon, lie within $0.01$ of full-precision Muon at equal wall-clock (Figure~\ref{fig:nanogpt}, Table~\ref{tab:nanogpt}). On the sign-after placement, error feedback is free: those two differ by $0.002$, below the noise level. The target they compress is an orthogonal factor whose entries stay evenly spread, the case in which the scaled sign loses least. The two models of EF21-MuonSign separate at this width. The run has a single client, so its broadcast model $\mathbf{W}$ is the only point at which a gradient is evaluated; there the compressed downlink costs nothing: $\mathbf{W}$ is indistinguishable from EF21-MuonUSign. The server model $\mathbf{X}$, which the guarantee bounds, settles $2.2$ nats above it. This is the mechanism of Remark~\ref{rem:dimension}, not a tuning failure, and Appendix~\ref{app:nanogpt} localizes it to the zero-initialized MLP output projections.

\section{Conclusion}\label{sec:conclusion}

\cl{Every placement of the sign around the Muon oracle can turn the update into an ascent direction on a linear objective, at every step size and momentum: SignMuon after the oracle, MuonUSign before it, MuonSign on both sides.} Whether error feedback repairs \cl{this} depends on where it is applied. Applied to the \emph{gradient}, it attains the standard nonconvex rate; applied to the \emph{polar factor}, it fails for every $(L,\eta,\mu)$. \cl{In all three settings the strongest compressed method signs \emph{after} the oracle, where no guarantee is available, and the two provably convergent methods trail it.} The counterexamples rule out a guarantee for every problem in the class. \cl{Whether a guarantee holds under further conditions that ordinary problems satisfy is the open question.}
